%% Motivation and Contextualization
\chapter*{Contextualization and Motivation}
Interventional radiology refers to image-guided therapeutic procedures
performed under real-time X-ray visualization \cite{6}. The number of these procedures has
increased over the past two decades worldwide. Interventional radiologists
and their teams are the most exposed group of medical radiation workers due to the prolonged
use of fluoroscopy and the proximity required during interventions \cite{10,55}.
Currently, protection strategies rely on lead garments
and structural shielding within the room \cite{37}. Compliance with established occupational dose limits is monitored
using passive dosimeters, which are worn either above or below the lead apron. These devices
are read \emph{a posteriori}, in the best-case scenario, once a month. As a result, there is no
immediate feedback during procedures and exposure incidents cannot be prevented \cite{46,67}.

Active personal dosimeters offer a complementary solution by providing real-time auditory and/or visual feedback. The main benefit
of these devices in IR is not to prevent exceedance of annual dose limits, which have been extensively studied and are rarely
surpassed under standard conditions. Instead, ADPs support
the principle of keeping exposure {\it as low as reasonably achievable }(ALARA) \cite{1,2,3}. These systems alert personnel to high-dose situations, promoting behavioural adjustments to reduce avoidable exposure. Adjustments may include reducing fluoroscopy time and frame rate, optimizing patient and operator positioning, and improve radiation shielding placement \cite{4}.

The proposed project explores the development of an APD system for IR, capable of monitoring
radiation exposure at four anatomical regions: eyes, neck, thorax (chest) and hands. The system will transmit dose data wirelessly to an external
dashboard for real-time visualization.
